\documentclass[11pt]{article}
\usepackage[T1]{fontenc}
\usepackage{textcomp}
\usepackage[margin=0.75in]{geometry}
\usepackage{amsmath,amssymb,amsthm}
\usepackage{array,booktabs}
\usepackage{hyperref}
\setlength{\parindent}{0pt}
\newtheorem{theorem}{Theorem}
\theoremstyle{definition}
\newtheorem{definition}{Definition}
\title{Flow Proof Artifact: prop-13-componendo-dividendo}
\author{Generated by \texttt{flow doc proof}}
\date{Flow Proof Book}
\begin{document}
\maketitle
If four magnitudes are proportional, they are proportional by composition and division.\\[0.5em]
\textbf{Source.} euclid — Elements, Book V, Proposition 13\\[0.5em]
\bigskip
\noindent{\large \textbf{Derived fact 1.} \textit{If four magnitudes are proportional, they are proportional by composition and division} \hfill the Euclidean plane \cdot Euclid Book V \cdot Proposition 13: proportional magnitudes satisfy componendo and dividendo}\par
\smallskip\noindent\textit{Coordinate: the Euclidean plane \cdot Euclid Book V \cdot Proposition 13: proportional magnitudes satisfy componendo and dividendo}\par
\smallskip\noindent\textit{Source: euclid — Elements, Book V, Proposition 13}\par
\smallskip\noindent\textit{Built on: proposition 11: proportional magnitudes satisfy componendo, for Euclid Book V on the Euclidean plane, proposition 12: proportional magnitudes satisfy dividendo, for Euclid Book V on the Euclidean plane}\par
\medskip
\noindent\textbf{Goal.} If four magnitudes are proportional, they are proportional by composition and division\par
\noindent$\displaystyle \text{ratio a plus b a minus b}$\par
\medskip
\noindent
\begin{tabular}{@{} >{\raggedright\arraybackslash}p{0.06\textwidth} >{\raggedright\arraybackslash}p{0.47\textwidth} >{\raggedleft\arraybackslash}p{0.41\textwidth} @{}}
\textbf{\#} & \textbf{Proof} & \textbf{Mathematics} \\
\midrule
\textbf{1.} & We prove that proposition 13: proportional magnitudes satisfy componendo and dividendo for Euclid Book V the Euclidean plane. & \\[0.45em]
\textbf{2.} & Let a = first. & \\[0.45em]
\textbf{3.} & Let b = second. & \\[0.45em]
\textbf{4.} & Let c = third. & \\[0.45em]
\textbf{5.} & Let d = fourth. & \\[0.45em]
\textbf{6.} & From step 2, step 3, step 4, and step 5, we can deduce that ratio(a plus b, a - b) equals ratio(c plus d, c - d). & $\displaystyle ratio(a + b, a - b) = ratio(c + d, c - d)$ \\[0.45em]
\textbf{7.} & We invoke the derived fact governing Euclid Book V the Euclidean plane: proposition 11: proportional magnitudes satisfy componendo, for Euclid Book V on the Euclidean plane. & \\[0.45em]
\textbf{8.} & We invoke the derived fact governing Euclid Book V the Euclidean plane: proposition 12: proportional magnitudes satisfy dividendo, for Euclid Book V on the Euclidean plane. & \\[0.45em]
\textbf{9.} & From step 7 and step 8, this implies ratio a plus b a minus b. Hence proven. & $\displaystyle \text{ratio a plus b a minus b}$ \\[0.45em]
\bottomrule
\end{tabular}
\label{thm:Geometry-Euclid Book V-proposition_13_proportional_magnitudes_satisfy_componendo_and_dividendo}
\par\medskip\hrule
\end{document}
